\documentclass[12pt,a4paper]{article}

% Packages
\usepackage[utf8]{inputenc}
\usepackage[spanish]{babel}
\usepackage{graphicx}
\usepackage{listings}
\usepackage{xcolor}
\usepackage{geometry}
\usepackage{float}

\geometry{margin=2.5cm}

\lstset{
    language=Prolog,
    basicstyle=\ttfamily\small,
    keywordstyle=\color{blue},
    commentstyle=\color{green!60!black},
    stringstyle=\color{red},
    showstringspaces=false,
    frame=single,
    breaklines=true
}

% portada 
\begin{document}

\begin{titlepage}
    \centering

    % Logo opcional
    %\includegraphics[width=0.3\textwidth]{logo.png}\par\vspace{1cm}

    {\scshape\Large Programación funcional \par}
    \vspace{1.5cm}

    {\huge\bfseries Reporte de Práctica en SWI-Prolog \par}
    \vspace{0.5cm}
    {\Large Tarea 2 \par}

    \vfill

    {\large
    Nombre del Alumno: Emilio Izquierdo Montero {\hspace{6cm}} \\
    \vspace{0.3cm}
    Profesor: Octavo Guzman Peñaloza {\hspace{6cm}} \\
    \vspace{0.3cm}
    Grupo / Semestre: 8 {\hspace{6cm}} \\
    \vspace{0.3cm}
    Fecha: 26 de mayo del 2026 {\hspace{6cm}}
    \par}

    \vfill

\end{titlepage}

\newpage

\section{Introducción}

La programación lógica es un paradigma de programación basado en la lógica formal y el razonamiento matemático, a diferencia de otros paradigmas donde el programador indica paso a paso cómo resolver un problema, en la programación lógica se describen hechos, relaciones y reglas que representan conocimiento.

Uno de los lenguajes más representativos de este paradigma es Prolog, el cual permite representar conocimiento y generar inferencias automáticas mediante consultas lógicas.

La diferencia principal entre el paradigma imperativo y el declarativo consiste en que los lenguajes imperativos indican cómo resolver un problema paso a paso, mientras que en los declarativos se describe qué se desea obtener.

La importancia de Prolog en Inteligencia Artificial radica en su capacidad para representar conocimiento simbólico, realizar inferencias automáticas y construir sistemas expertos, motores de recomendación y aplicaciones basadas en razonamiento lógico.

\section{Base de Conocimientos}

\subsection{Hechos y Reglas sobre Startups}

\begin{lstlisting}

% hechos

startup(innovaciondudosa).
startup(leppe).
startup(lotus).
startup(rebonet).
startup(luzmaya).
startup(smartguide).
startup(laika).
startup(centinela).

usa_tecnologia(innovaciondudosa, agente_IA).
usa_tecnologia(innovaciondudosa, base_de_datos).

usa_tecnologia(leppe, tecnologia_emergente).
usa_tecnologia(leppe, inteligencia_artificial).
usa_tecnologia(leppe, vision_computacional).

usa_tecnologia(lotus, redes_neuronales).
usa_tecnologia(lotus, vision_computacional).
usa_tecnologia(lotus, inteligencia_artificial).

usa_tecnologia(rebonet, industria_pecuaria).
usa_tecnologia(rebonet, material_organico).

usa_tecnologia(luzmaya, energias_renovables).

usa_tecnologia(smartguide, iot).
usa_tecnologia(smartguide, industria40).

usa_tecnologia(laika, redes_neuronales).
usa_tecnologia(laika, vision_computacional).

usa_tecnologia(centinela, redes_neuronales).
usa_tecnologia(centinela, vision_computacional).

especialistas(luis, vision_computacional).
especialistas(luis, inteligencia_artificial).

especialistas(emilio, vision_computacional).
especialistas(emilio, inteligencia_artificial).
especialistas(emilio, redes_neuronales).
especialistas(emilio, repositorios).

especialistas(brian, inteligencia_artificial).

especialistas(teresa, iot).
especialistas(teresa, industria40).

especialistas(rene, energias_renovables).
especialistas(rene, iot).

inversion_economica(leppe).
inversion_economica(rebonet).

% reglas

startup_ia(X):-
    usa_tecnologia(X, inteligencia_artificial).

startup_exitosa(X):-
    usa_tecnologia(X, inteligencia_artificial),
    inversion_economica(X).

startup_sustentable(X):-
    usa_tecnologia(X, energias_renovables);
    usa_tecnologia(X, material_organico).

startup_futurista(X):-
    usa_tecnologia(X, inteligencia_artificial),
    usa_tecnologia(X, vision_computacional).

colaborador_ideal(Persona, Startup):-
    especialistas(Persona, Tecnologia),
    usa_tecnologia(Startup, Tecnologia).

ecosistema_avanzado(X):-
    usa_tecnologia(X, redes_neuronales),
    usa_tecnologia(X, vision_computacional).

\end{lstlisting}

\subsection{Hechos y Reglas sobre Pensadores}

\begin{lstlisting}


filosofo(descartes).
filosofo(platon).
filosofo(aristoteles).
filosofo(socrates).

fisico(einstein).
fisico(newton).
fisico(galileo).
fisico(tesla).

matematico(newton).
matematico(einstein).
matematico(gauss).
matematico(euler).

quimico(bohr).
quimico(curie).
quimico(lavoisier).
quimico(dalton).

astronomo(copernico).
astronomo(hubble).
astronomo(webb).
astronomo(galileo).


maestro(einstein, tesla).
maestro(lavoisier, curie).
maestro(copernico, galileo).
maestro(platon, aristoteles).


sabe_matematicas(X) :-
    matematico(X).

sabe_astronomia(X) :-
    astronomo(X).

sabe_quimica(X) :-
    quimico(X).

sabe_filosofia(X) :-
    filosofo(X).

sabe_fisica(X) :-
    fisico(X).


cientifico_integral(X) :-
    fisico(X),
    matematico(X).

pensador(X) :-
    filosofo(X),
    fisico(X).

erudito(X) :-
    matematico(X),
    filosofo(X).

erudito(X) :-
    fisico(X),
    filosofo(X).

\end{lstlisting}

\section{Evidencias de Ejecución}

\subsection{Carga del Archivo en SWI-Prolog}

\begin{figure}[H]
    \centering
    \includegraphics[width=0.9\linewidth]{/home/emilio/Descargas/cargapro.png}
    \caption{Carga correcta del archivo .pro}
\end{figure}

\subsection{Consulta con Variables}

\begin{figure}[H]
    \centering
    \includegraphics[width=0.9\linewidth]{/home/emilio/Descargas/variables.png}
    \caption{Consulta utilizando variables}
\end{figure}

%-----------------------------------------------------

\subsection{Uso de Reglas}

\begin{figure}[H]
    \centering
    \includegraphics[width=0.9\linewidth]{/home/emilio/Descargas/reclas.png}
    \caption{Uso de reglas lógicas}
\end{figure}

%-----------------------------------------------------

\subsection{Backtracking}

\begin{figure}[H]
    \centering
    \includegraphics[width=0.9\linewidth]{/home/emilio/Descargas/final.png}
    \caption{Resultados múltiples mediante backtracking}
\end{figure}

%=====================================================
\section{Explicación de Resultados}
%=====================================================

Durante la ejecución de las consultas, Prolog utilizó los hechos y reglas almacenados en la base de conocimiento para responder automáticamente a cada pregunta realizada.

En las consultas simples, el sistema únicamente verificó si el hecho existía dentro de la base de conocimientos.

En las consultas con variables, Prolog realizó un proceso de búsqueda para encontrar todos los valores que cumplían las condiciones indicadas.

Las reglas permitieron generar nuevo conocimiento a partir de hechos ya existentes. Por ejemplo, mediante la regla \texttt{startup\_exitosa(X)}, el sistema pudo inferir qué empresas utilizaban inteligencia artificial y además contaban con inversión económica.

El proceso de inferencia consiste en deducir nueva información utilizando reglas lógicas. Asimismo, el backtracking permitió explorar múltiples soluciones posibles cuando existía más de una respuesta válida para una consulta.

\section{Análisis Técnico y Reflexión}

La programación lógica presenta diferencias importantes respecto a lenguajes imperativos como Java, Python o C++. En los lenguajes tradicionales se especifica paso a paso cómo resolver un problema mediante instrucciones y estructuras de control.

En contraste, Prolog trabaja describiendo hechos y relaciones, permitiendo que el motor lógico determine automáticamente las respuestas mediante inferencia.

Entre las ventajas observadas se encuentran la facilidad para representar conocimiento, la capacidad de generar inferencias y la implementación sencilla de sistemas expertos.

Una de las principales dificultades fue comprender el funcionamiento del backtracking y la forma en que Prolog evalúa las reglas para encontrar soluciones.

Este paradigma puede aplicarse en sistemas inteligentes reales como motores de recomendación, asistentes virtuales, sistemas expertos médicos, análisis semántico y aplicaciones de Inteligencia Artificial simbólica.

\section{Conclusiones}

La práctica permitió comprender el funcionamiento básico de la programación lógica utilizando Prolog. Se observó cómo es posible representar conocimiento mediante hechos y reglas, además de generar inferencias automáticas utilizando consultas lógicas.

\end{document}
